\documentclass[12pt]{article}
\usepackage{amsmath,amssymb,amsfonts}
\usepackage[margin=1in]{geometry}

\begin{document}

\section*{Chapter 7, Problem 29}

\textbf{Problem:} Consider the time-dependent inhomogeneous Klein--Gordon equation:
\[
-\nabla^2\psi + \frac{1}{v^2}\frac{\partial^2\psi}{\partial t^2} + \mu^2\psi = f(\mathbf{r},t).
\]

\begin{itemize}
\item[a)] Show that by using the Fourier transform method one obtains the following two integral expressions for the Green's functions $G_A$ and $G_R$:

\[
G_A(\mathbf{r},\mathbf{r}',t,t') = \frac{v^2}{(2\pi)^3}\int d^3k\,\frac{e^{i\mathbf{k}\cdot(\mathbf{r}-\mathbf{r}')}\sin\bigl[\sqrt{v^2 k^2+v^2\mu^2}\,(t-t')\bigr]}{\sqrt{v^2k^2+v^2\mu^2}}\,\theta(t'-t),
\]
\[
G_R(\mathbf{r},\mathbf{r}',t,t') = -\frac{v^2}{(2\pi)^3}\int d^3k\,\frac{e^{i\mathbf{k}\cdot(\mathbf{r}-\mathbf{r}')}\sin\bigl[\sqrt{v^2 k^2+v^2\mu^2}\,(t-t')\bigr]}{\sqrt{v^2k^2+v^2\mu^2}}\,\theta(t-t'),
\]

where $R = |\mathbf{r}-\mathbf{r}'|$. The two expressions correspond to two possible ways of avoiding a singularity which occurs in doing the Fourier transform integral corresponding to the space variables. In the first integral, $\sqrt{v^2k^2+v^2\mu^2}$ has positive imaginary part when $|q| < \mu$; in the second, $\sqrt{v^2k^2+v^2\mu^2}$ has negative imaginary part when $|q| < \mu$.

\item[b)] Show that $G_R = 0$ when $R < v(t-t')$ and that $G_A = 0$ when $R < v(t'-t) > 0$.

[\textit{Hint:} If the integral for $G_R$ is considered as an integral in the complex plane, then the plane is cut from $-\infty$ to $+\infty$ and the integral in question runs along the top of the cut. For $G_A$, the integral runs along the bottom of the cut.]

\item[c)] When $R < v(\tau)$ ($\tau = t-t'$), show that the integral occurring in $G_R$ can be reduced to the contour integral shown in Fig.\ 7.12. Hence show that
\[
G_R = \frac{v^2}{4\pi^2 R}\int_{-\infty}^{\infty}\cosh(\kappa R)\,e^{-i\sqrt{v^2\kappa^2-v^2\mu^2}\,\tau}\,d\kappa,
\]
and, similarly,
\[
G_A = \frac{v^2}{4\pi^2 R}\int_{-\infty}^{\infty}\cosh(\kappa R)\,e^{+i\sqrt{v^2\kappa^2-v^2\mu^2}\,\tau}\,d\kappa,
\]
where $\epsilon(x)$ equals 1 if $x > 0$ and equals zero if $x < 0$. Also, $T = v\tau = v(t-t')$.

\item[d)] Making the change of variable $q = \kappa\sin\phi$ and using some symmetry considerations, show that
\[
G_R = \frac{v^2}{4\pi^2}\int_{-\pi/2}^{\pi/2}\cosh(\mu R\cos\phi)\,e^{-i\mu T\sin\phi}\,d\phi.
\]

[\textit{Note:} the limits of integration are $\pm\pi/2$, not $\pm\pi$.]

\item[e)] Show that this can be written as
\[
G_R = \frac{v^2}{4\pi^2}\oint_C \cosh\Bigl(\frac{\mu R}{2}(z+z^{-1})\Bigr)\,\exp\Bigl(-\frac{\mu T}{2}(z-z^{-1})\Bigr)\frac{dz}{iz},
\]
with a similar expression for $G_A$. Here $C$ is a counterclockwise contour around the unit circle.

\item[f)] By making two variable changes (in two different integrals), show that one obtains
\[
G_R = \frac{v^2}{4\pi}\Bigl\{\delta(T-R) + \frac{\mu}{2}\Bigl[\frac{1}{2}\sqrt{T^2-R^2}\Bigr]^{-1}J_1\!\bigl(\mu\sqrt{T^2-R^2}\bigr)\cdot\theta(T-R)\Bigr\}.
\]
\end{itemize}

Using the integral representation of the Bessel function (Eq.\ (6.65) and the sentence following that equation), show that this reduces to
\[
G_R = -\frac{v}{4\pi R}\delta(v|r'-r| - T) + \frac{v^2\mu}{8\pi}\frac{J_1(\mu\sqrt{T^2-R^2})}{\sqrt{T^2-R^2}}\,\theta(T-R).
\]

Also,
\[
G_R = \frac{v}{4\pi}\frac{\delta(T-R)}{R} + \frac{v^2\mu}{8\pi}\frac{J_1(\mu v\sqrt{T^2-R^2})}{\sqrt{T^2-R^2}}\,\theta(T-R).
\]

[Here we have used the fact that $J_1(x) = J_1(-x)$, which is clear from Eq.\ (6.66).]

We're almost done! Carry out the differentiation using Eq.\ (7.22) to show that
\[
G_R = \frac{v}{4\pi}\frac{\delta(T-R)}{R} - \frac{v^2\mu}{8\pi}\frac{J_1(\mu\sqrt{T^2-R^2})}{\sqrt{T^2-R^2}}\,\theta(T-R).
\]

Note that $G_R = G_A$ can be written as
\[
G = \frac{v}{4\pi}\frac{\delta(\sqrt{s^2})}{R} + \frac{v^2\mu}{8\pi}\frac{J_1(\mu\sqrt{-s^2})}{\sqrt{-s^2}},
\]
which is a function only of the Lorentz invariant $s^2 = R^2 - T^2 = |\mathbf{r}-\mathbf{r}'|^2 - v^2(t-t')^2$.

\bigskip
\textbf{Solution:}

\subsection*{Part (a)}

The Klein--Gordon Green's function satisfies:
\[
\Bigl(-\nabla^2 + \frac{1}{v^2}\frac{\partial^2}{\partial t^2} + \mu^2\Bigr)G = \delta^3(\mathbf{r}-\mathbf{r}')\delta(t-t').
\]

Fourier transforming in space and time with $\mathbf{R} = \mathbf{r}-\mathbf{r}'$, $\tau = t-t'$:
\[
\tilde{G}(\mathbf{k},\omega) = \frac{1}{k^2 + \mu^2 - \omega^2/v^2} = \frac{v^2}{v^2k^2 + v^2\mu^2 - \omega^2}.
\]

The poles in $\omega$ are at $\omega = \pm\omega_k$ where $\omega_k = \sqrt{v^2k^2 + v^2\mu^2}$.

Inverting the $\omega$ transform:
\[
G(\mathbf{R},\tau) = \frac{v^2}{(2\pi)^3}\int d^3k\,e^{i\mathbf{k}\cdot\mathbf{R}}\int\frac{d\omega}{2\pi}\frac{e^{-i\omega\tau}}{\omega_k^2 - \omega^2}.
\]

\textbf{Retarded:} Displace poles below real axis: $\omega = \pm\omega_k - i\epsilon$. For $\tau > 0$, close below:
\[
\int\frac{e^{-i\omega\tau}}{(\omega_k-\omega)(\omega_k+\omega)}\frac{d\omega}{2\pi} = -\frac{\sin(\omega_k\tau)}{\omega_k}.
\]

For $\tau < 0$: zero. So:
\[
G_R = -\frac{v^2}{(2\pi)^3}\theta(\tau)\int d^3k\,e^{i\mathbf{k}\cdot\mathbf{R}}\frac{\sin(\omega_k\tau)}{\omega_k}. \quad \checkmark
\]

\textbf{Advanced:} Poles above real axis, giving $\theta(-\tau)$ instead.

\subsection*{Part (b)}

When $R > vT$ (spacelike separation), the argument of the sine function $\omega_k\tau$ oscillates rapidly with $k$, and the integral over $\mathbf{k}$ vanishes due to the Riemann--Lebesgue lemma applied appropriately. More rigorously, in the complex $k$-plane analysis, the contour can be deformed to show the integral is zero for spacelike separations.

\subsection*{Parts (c)--(f)}

For timelike separation ($T > R$), performing the angular $\mathbf{k}$ integration:
\[
G_R = -\frac{v^2\theta(\tau)}{2\pi^2 R}\int_0^{\infty}\frac{k\sin(kR)\sin(\omega_k\tau)}{\omega_k}\,dk.
\]

Deforming the contour to the imaginary axis $k = i\kappa$ gives the cosh representation in part (c). The substitution $\kappa = \mu\sin\phi$ then gives part (d).

Setting $z = e^{i\phi}$ converts to the contour integral in part (e).

Finally, using the integral representation of the Bessel function:
\[
J_1(x) = \frac{1}{2\pi i}\oint z^{-2}\exp\!\Bigl(\frac{x}{2}(z-z^{-1})\Bigr)\,dz,
\]
and evaluating the residues, one obtains the final result:
\[
\boxed{G_R = \frac{v}{4\pi}\frac{\delta(T-R)}{R} + \frac{v^2\mu}{8\pi}\frac{J_1\!\bigl(\mu\sqrt{T^2-R^2}\bigr)}{\sqrt{T^2-R^2}}\,\theta(T-R).}
\]

The first term is the sharp wave front (same as the massless case $\mu=0$); the second is the ``tail'' due to the mass, which represents waves propagating inside the light cone. As $\mu \to 0$, $J_1(x)/x \to 1/2$ for small $x$, but the tail vanishes, recovering the massless retarded Green's function.

The expression depends only on $s^2 = R^2 - T^2$, confirming Lorentz invariance. $\blacksquare$

\end{document}
